\documentclass[a4paper,11pt]{article}
\usepackage[utf8]{inputenc}
\usepackage[T1]{fontenc}
\usepackage{lmodern}
\usepackage{amsmath,amssymb,amsthm}
\usepackage{mathtools}
\usepackage{booktabs}
\usepackage{longtable}
\usepackage{array}
\usepackage{hyperref}
\usepackage{graphicx}
\usepackage{float}
\usepackage{geometry}
\geometry{margin=25mm}

\hypersetup{
  colorlinks=true,
  linkcolor=blue,
  urlcolor=blue,
  citecolor=blue
}

\theoremstyle{definition}
\newtheorem{theorem}{Theorem}[section]
\newtheorem{proposition}[theorem]{Proposition}
\newtheorem{lemma}[theorem]{Lemma}
\newtheorem{corollary}[theorem]{Corollary}
\newtheorem{definition}[theorem]{Definition}
\newtheorem{remark}[theorem]{Remark}
\newtheorem{observation}[theorem]{Observation}

\setlength{\parindent}{1em}
\setlength{\parskip}{0.5em}

\providecommand{\tightlist}{\setlength{\itemsep}{0pt}\setlength{\parskip}{0pt}}
\begin{document}

\section{On the Limits of the Curvature Radius in Central Projection}\label{on-the-limits-of-the-curvature-radius-in-central-projection}

\textbf{The Cases \(R \to \infty\) and \(R \to 0\)}

\begin{center}\rule{0.5\linewidth}{0.5pt}\end{center}

\textbf{Author}: Noriaki Kihara (WF System Co., Ltd.)

\textbf{Date}: April 2026

\textbf{DOI:} \href{https://doi.org/10.5281/zenodo.19526549}{10.5281/zenodo.19526549}

\begin{center}\rule{0.5\linewidth}{0.5pt}\end{center}

\subsection{Abstract}\label{abstract}

This paper examines both limits \(R \to \infty\) and \(R \to 0\) of the curvature radius \(R\) within the central projection framework established in the preceding papers {[}1{]}, {[}2{]}, {[}3{]}. In both limits, the induced metric, the curvature tensor, and each term of the derived Einstein equations are geometrically regular. In the \(R \to 0\) limit, however, a structural tension arises in relation to the upper bound on geodesic deviation established in paper {[}2{]}. As one method for resolving this tension within the geometric framework itself, we show that a \textbf{nested structure of subjective spaces} can be constructed directly from the existing geometric degrees of freedom in the preceding papers {[}1{]}, {[}2{]}. This paper asserts only the geometric existence of such a structure and makes no claim regarding its physical reality.

\begin{center}\rule{0.5\linewidth}{0.5pt}\end{center}

\subsection{1. Introduction}\label{introduction}

In the preceding paper {[}1{]}, the central projection

\[
\Phi : \Pi_R \longrightarrow S^n(R)
\]

between the hypersphere \(S^n(R)\) of radius \(R\) in the \((n+1)\)-dimensional Euclidean background space \(\mathbb{R}^{n+1}\) and the tangent hyperplane \(\Pi_R\) was defined, and it was shown that the metric induced on \(\Pi_R\),

\[
g_{\mu\nu}(x) = \frac{R^2}{l^2}\left(\delta_{\mu\nu} - \frac{x_\mu x_\nu}{l^2}\right), \qquad l^2 = R^2 + |x|^2
\]

yields the Einstein equations

\[
G_{\mu\nu} + \Lambda\, g_{\mu\nu} = 0, \qquad \Lambda = \frac{3}{R^2}.
\]

In paper {[}2{]}, the symmetries of this framework (discrete stability, axis equivalence, the structure of geodesic deviation, and the transformability of subjective coordinate systems) were proved. In particular, it was shown that the geodesic deviation \(|\xi|^2 = g_{\mu\nu}\xi^\mu \xi^\nu\) measured by an observer within the subjective space has an upper bound depending on the curvature radius \(R\).

In paper {[}3{]}, the relativity of observation when multiple subjective spaces share the same background space was discussed.

In these papers, the curvature radius \(R\) remains a free parameter, and the construction is well-defined throughout \(R \in (0, \infty)\). The purpose of this paper is to examine the geometric structure at both limiting cases

\[
R \to \infty, \qquad R \to 0
\]

and to record the questions that naturally arise from this examination.

\begin{center}\rule{0.5\linewidth}{0.5pt}\end{center}

\subsection{\texorpdfstring{2. The Case \(R \to \infty\)}{2. The Case R \textbackslash to \textbackslash infty}}\label{the-case-r-to-infty}

\subsubsection{2.1 Geometric Regularity}\label{geometric-regularity}

In the limit \(R \to \infty\), the induced metric \(g_{\mu\nu}\) asymptotically approaches

\[
g_{\mu\nu}(x) \longrightarrow \delta_{\mu\nu}
\]

and the subjective space \(\Pi_R\) approaches a flat Euclidean space. The cosmological term becomes

\[
\Lambda = \frac{3}{R^2} \longrightarrow 0,
\]

and each term of the Einstein equations exhibits a mathematically well-defined asymptotic behavior. Each component of the curvature tensor approaches zero continuously, and the geometric construction of paper {[}1{]} does not break down in this limit.

\subsubsection{2.2 Questions That Remain Outside the Geometry}\label{questions-that-remain-outside-the-geometry}

However, when one attempts a physical interpretation of this construction, the following questions arise. These are questions that lie outside the geometric framework and are not addressed in this paper.

\textbf{Question 2.1}　The background space \(\mathbb{R}^{n+1}\) carries a Euclidean signature. Whether the axis corresponding to time is structurally distinguished when the subjective space becomes flat in the limit \(R \to \infty\) is not determined within the framework of pure geometry.

\textbf{Question 2.2}　While \(\Lambda \to 0\) is a smooth asymptotic, the volume of the subjective space diverges as \(R \to \infty\). The behavior of quantities that may appear as products of these (quantities that could be interpreted as vacuum energy) depends on interpretations defined outside the geometry.

\textbf{Question 2.3}　When \(\Lambda\) is derived from a framework other than central projection, whether its behavior is consistent with \(3/R^2\) in the limit \(R \to \infty\) depends on the structure of each respective framework.

\textbf{Key point}: The construction is geometrically regular, but divergence in physical interpretation cannot be excluded.

\begin{center}\rule{0.5\linewidth}{0.5pt}\end{center}

\subsection{\texorpdfstring{3. The Case \(R \to 0\)}{3. The Case R \textbackslash to 0}}\label{the-case-r-to-0}

\subsubsection{3.1 Geometric Regularity}\label{geometric-regularity-1}

In the limit \(R \to 0\), \(\Lambda = 3/R^2\) diverges. However, the induced metric \(g_{\mu\nu}\), the curvature tensor, and each term of the Einstein equations are well-defined for any \(R > 0\), and the limit \(R \to 0\) is also defined as a geometrically regular asymptotic. The construction of paper {[}1{]} does not break down in this limit either.

\subsubsection{3.2 Structural Tension with the Upper Bound on Geodesic Deviation}\label{structural-tension-with-the-upper-bound-on-geodesic-deviation}

In paper {[}2{]}, it was shown that the geodesic deviation \(|\xi|\) measured by an observer within the subjective space has an upper bound depending on the curvature radius \(R\). That is, the quantities measurable by the observer are constrained by the order of magnitude of the subjective space.

In the limit \(R \to 0\), this upper bound asymptotes to zero while the local curvature diverges at the order of \(1/R^2\). A structural tension therefore arises between the quantities measurable by an observer within the subjective space and the local geometric structure of the subjective space in which that observer resides.

\subsubsection{3.3 Questions That Remain Outside the Geometry}\label{questions-that-remain-outside-the-geometry-1}

\textbf{Question 3.1}　\(\Lambda = 3/R^2 \to \infty\) is defined as the divergence of a geometric quantity, but how this divergence is to be interpreted physically is a question outside the geometry.

\textbf{Question 3.2}　In the limit where the subjective space contracts to a point, whether the structure itself --- that an observer exists within the interior of the subjective space --- is preserved is not determined within the geometry.

\textbf{Key point}: The construction is geometrically regular, but divergence in physical interpretation cannot be excluded, and the tension noted in §3.2 leaves open questions regarding the consistency with the observer's structure.

\begin{center}\rule{0.5\linewidth}{0.5pt}\end{center}

\subsection{4. Resolution of the Tension: Nested Structure of Subjective Spaces}\label{resolution-of-the-tension-nested-structure-of-subjective-spaces}

We present one method for resolving the tension identified in §3.2 within the framework of the geometry itself. What is described in this section is a direct consequence of existing theorems in the preceding papers {[}1{]}, {[}2{]}, and requires no introduction of new axioms.

\subsubsection{4.1 Construction of the Nested Structure}\label{construction-of-the-nested-structure}

In the central projection framework of paper {[}1{]}, the subjective space \(\Pi_R\) is defined by the choice of axis \(e\). By Theorem 3.1 of paper {[}2{]} (axis equivalence), this choice of axis is not geometrically unique, and there is a degree of freedom to choose any axis \(e'\) as a new projection center.

This degree of freedom is applicable to any point in the interior of the subjective space. That is, for any point \(p \in \Pi_R\), one can define another central projection

\[
\Phi' : \Pi_{R'} \longrightarrow S^n(R')
\]

with \(p\) as the new base point of the axis. The new curvature radius \(R' > 0\) can be chosen freely. The new subjective space \(\Pi_{R'}\) has the same geometric structure as the construction in paper {[}1{]} (induced metric, curvature tensor, Einstein equations).

\subsubsection{4.2 Existence Proposition}\label{existence-proposition}

\textbf{Proposition 4.1} (Existence of Nested Structure)　In the central projection framework defined in paper {[}1{]}, for any subjective space \(\Pi_R\) and any point \(p \in \Pi_R\) in its interior, a subjective space \(\Pi_{R'}\) exists for any \(R' > 0\) satisfying the following properties:

\begin{enumerate}
\def\labelenumi{\arabic{enumi}.}
\tightlist
\item
  \(\Pi_{R'}\) is defined by a new central projection centered at \(p\).
\item
  \(\Pi_{R'}\) has the same geometric structure as the construction defined in paper {[}1{]}.
\item
  \(\Pi_{R'}\) is constructed from the same background space \(\mathbb{R}^{n+1}\) as \(\Pi_R\).
\end{enumerate}

\textbf{Proof Sketch}　By Theorem 3.1 of paper {[}2{]}, at any interior point \(p\) of \(\Pi_R\), an axis \(e'\) with \(p\) as its base point can be chosen. From this \(e'\), the hypersphere \(S^n(R')\) of radius \(R'\) and the tangent hyperplane \(\Pi_{R'}\) can be constructed by the same procedure as in paper {[}1{]}. The construction is well-defined by the theorems of paper {[}1{]}. \(\square\)

\subsubsection{4.3 Iterability of the Nesting}\label{iterability-of-the-nesting}

The construction of Proposition 4.1 is iterable. At an interior point \(p'\) of \(\Pi_{R'}\), one can further raise \(\Pi_{R''}\), and so on. This iteration is closed within the framework of papers {[}1{]}, {[}2{]} and requires no additional axioms.

\subsubsection{4.4 Relation to the Tension in §3.2}\label{relation-to-the-tension-in-3.2}

The existence of Proposition 4.1 allows the tension identified in §3.2 to be organized within the geometry as follows:

For a region of the subjective space \(\Pi_R\) that is in the limiting regime \(R \to 0\), there exists a geometric degree of freedom to interpret that region as the base point of a new subjective space \(\Pi_{R'}\). A region that, viewed from the parent subjective space \(\Pi_R\), appears to be degenerating, has a well-defined geometric structure when viewed from the interior of the new central projection \(\Pi_{R'}\).

Through this reorganization, structural elements that become difficult to accommodate within the parent subjective space have a geometric degree of freedom to be described as a projection onto the child subjective space.

\begin{center}\rule{0.5\linewidth}{0.5pt}\end{center}

\subsection{5. Discussion}\label{discussion}

In this paper, we examined both limits of the curvature radius of the central projection framework of the preceding papers {[}1{]}, {[}2{]}, {[}3{]}, and showed in §4 that the nested structure of subjective spaces can be constructed as a geometric degree of freedom within the existing framework.

The claims of this paper are limited to the following:

\begin{enumerate}
\def\labelenumi{\arabic{enumi}.}
\tightlist
\item
  In both limits, the geometric construction does not break down.
\item
  In the limit \(R \to 0\), the structural tension of §3.2 arises.
\item
  As a method for organizing this tension within the geometry, the existence proposition for the nested structure in §4 holds.
\end{enumerate}

This paper does not claim that such a nested structure physically exists. The questions presented in §2.2 and §3.3 also involve physical interpretations that are outside the geometry and are not addressed here.

Answers to these questions, and the attribution of physical meaning to the nested structure, are left to the judgment of the reader.

\begin{center}\rule{0.5\linewidth}{0.5pt}\end{center}

\subsection{6. Conclusion}\label{conclusion}

In the central projection framework of the preceding papers {[}1{]}, {[}2{]}, {[}3{]}, both limits \(R \to \infty\) and \(R \to 0\) of the curvature radius \(R\) are geometrically regular. The structural tension with the upper bound on geodesic deviation that arises in the limit \(R \to 0\) can be organized within the geometry by means of a nested structure of subjective spaces constructed from the existing degrees of freedom within the same framework. The physical interpretation of this construction lies outside the scope of this paper and is left to the reader.

\begin{center}\rule{0.5\linewidth}{0.5pt}\end{center}

\subsection{References}\label{references}

{[}1{]} Noriaki Kihara, ``Geometric Formulation via Central Projection,'' Zenodo, DOI: 10.5281/zenodo.19427780 (2025).

{[}2{]} Noriaki Kihara, ``Symmetries in the Central Projection Framework,'' Zenodo, DOI: 10.5281/zenodo.19434932 (2025).

{[}3{]} Noriaki Kihara, ``Relativity of Observation in Multiple Subjective Spaces,'' Zenodo, DOI: 10.5281/zenodo.19435162 (2025).

\begin{center}\rule{0.5\linewidth}{0.5pt}\end{center}

\end{document}
